% ============================================================
% SCHEDE PG AGGIUNTIVE — CASA FIESCHI
% Personaggi per tavolo alternativo / quinta sessione
% ============================================================

\chapter{Schede PG Aggiuntive --- Casa Fieschi}
\index{schede PG!Fieschi!aggiuntivi}

% ════════════════════════════════════════════════════════════
% PG 6 — VIOLANTE ADORNO
% @rpg-schema: <https://rpg-schema.org/instances/genova1507/pg-violante-adorno>
%   a fitd:Character ;
%   rdfs:label "Violante Adorno" ;
%   fitd:crew <.../crew-fieschi> ;
%   fitd:role "La Spia in Casa Francese / L'Infiltrata" ;
%   fitd:action [ fitd:actionName "Survey" ; fitd:dots 4 ] ;
%   fitd:action [ fitd:actionName "Prowl" ; fitd:dots 3 ] ;
%   fitd:action [ fitd:actionName "Sway" ; fitd:dots 3 ] .
% ════════════════════════════════════════════════════════════

\pgheader[FieschiGreen]{Violante Adorno}{La Spia in Casa Francese / L'Infiltrata}{Fieschi}
\pgquote{Lavoro per i Fieschi da sei anni. I francesi pensano che lavori per loro da cinque. Entrambe le cose sono vere.}

\section*{Background \& Motivazione}

\begin{rulebox}{}
  \textbf{Background:} Ventotto anni, di famiglia genovese minore ma ben collocata.
  È entrata al servizio dell'amministrazione francese come traduttrice e scrivana
  cinque anni fa — su indicazione dei Fieschi, che l'avevano reclutata l'anno prima.
  Da allora vive una doppia esistenza con un'efficienza che spaventa chiunque la
  conosca bene. Ha accesso agli archivi del governatore, conosce le comunicazioni
  cifrate dei francesi, e sa dove sono dislocate le truppe di riserva che
  Trivulzio non ha ancora mostrato.

  \medskip\noindent\textcolor{FieschiGreen}{\textbf{Motivazione:}} È stanca. Sei
  anni di doppio gioco logorano. Vuole che questa notte finisca con i francesi
  fuori da Genova in modo da poter essere un'unica persona invece di due.
\end{rulebox}

\section*{Attributi \& Azioni}


\begin{attributoblock}[FieschiGreen]{INSIGHT — Mente}
  \azione{Caccia}{Hunt}{2}{Trova informazioni negli archivi e nei dispacci}
  \azione{Studio}{Study}{3}{Lingue (francese, latino, catalano), codici militari}
  \azione{Osserva}{Survey}{4}{Niente le sfugge — è sopravvissuta grazie a questo}
  \azione{Ingegno}{Tinker}{2}{Codici cifrati, documenti falsi, sigilli}
\end{attributoblock}


\begin{attributoblock}[FieschiGreen]{PROWESS — Corpo}
  \azione{Finezza}{Finesse}{2}{Mani veloci, movimenti controllati}
  \azione{Ombra}{Prowl}{3}{Sa muoversi nel palazzo del governatore di notte}
  \azione{Combattimento}{Skirmish}{1}{Stiletto nascosto — ultima risorsa}
  \azione{Forza bruta}{Wreck}{0}{No}
\end{attributoblock}


\begin{attributoblock}[FieschiGreen]{RESOLVE — Spirito}
      \azione{Percezione}{Attune}{2}{Sente quando una copertura è a rischio}
    \azione{Comando}{Command}{1}{Autorità amministrativa — funziona sui funzionari}
    \azione{Relazioni}{Consort}{2}{Impiegati francesi, notai, interpreti}
    \azione{Persuasione}{Sway}{3}{Convince mentendo con la stessa facilità con cui dice la verità}
  \end{attributoblock}

\medskip
\begin{pgclockbox}{Stress \hfill \clock{9}{0}}
  \small\textit{Traumi: $\square$ Freddo \quad $\square$ Duro \quad $\square$ Ossessionato \quad $\square$ Paranoico \quad $\square$ Irrazionale \quad $\square$ Spezzato}
\end{pgclockbox}

\section*{Abilità Speciale}

% @rpg-schema: <.../pg-violante-adorno>
%   fitd:specialAbility [ rdfs:label "Accesso agli Archivi" ] .
\begin{fieschibox}{Accesso agli Archivi}
  Violante ha libero accesso all'ufficio del governatore e ai dispacci militari
  francesi. Una volta per sessione può dichiarare di conoscere un'informazione
  riservata dell'amministrazione occupante: posizione delle truppe, comunicazioni
  intercettate, ordini imminenti.

  \medskip\noindent\textit{Nessun tiro. L'informazione è accurata.
  Usarla operativamente richiede che non venga tracciata a lei —
  altrimenti la sua copertura salta.}
\end{fieschibox}

\section*{Legami}

\begin{table}[h]
\centering\renewcommand{\arraystretch}{1.4}
\begin{tabularx}{\linewidth}{@{}L{2.6cm} X@{}}
  \toprule\rowcolor{FieschiGreen}
  \textcolor{white}{\textbf{PG}} & \textcolor{white}{\textbf{Legame}} \\
  \midrule
  Niccolò & Il conte mi ha reclutata. Gli sono fedele. Non gli dico tutto perché alcune cose lo metterebbero in pericolo. \\
  \rowcolor{FieschiGreen!8}
  Frà Tomaso & Il predicatore sa che esisto ma non sa cosa faccio. Ogni tanto mi guarda come se stesse per capire. \\
  Ser Elia & Il capitano mi tratta come un soldato. È la cosa più rispettosa che qualcuno mi abbia mai fatto. \\
  \rowcolor{FieschiGreen!8}
  Dottor Ansaldo & Il medico. L'unico che sa dove mi nascondo quando le cose si mettono male. \\
  \bottomrule
\end{tabularx}
\end{table}

\section*{Equipaggiamento}
\begin{goldlist}
  \item Tessera di accesso all'ufficio del governatore — ancora valida
  \item Copia dei dispacci militari degli ultimi dieci giorni, in codice
  \item Cambio di abiti completo da impiegata francese e da popolana genovese
  \item Stiletto ceramico nascosto nell'orlo della gonna
\end{goldlist}


\begin{secretbox}
  \textbf{Segreto:} Violante ha scoperto che Monsignore Gregorio ha promesso
  informazioni ai francesi in cambio del supporto per l'arcivescovado.
  Lo sa da tre settimane. Non lo ha detto ai Fieschi perché Gregorio è
  il cugino del conte — e perché le prove che ha sono state ottenute
  tramite la sua copertura, che salterebbe se le usasse.
  Ha anche scoperto che i francesi sanno della sua doppia identità —
  ma hanno scelto di non agire ancora, perché la usano come canale
  per disinformare i Fieschi. Violante spia i francesi che la usano
  per spiare i Fieschi. Non è sicura di chi stia vincendo.

  \medskip\noindent\textcolor{SecretRed}{\textbf{Agenda:}}
  \textit{Rivelare la cosa su Gregorio questa notte — ma solo nel momento
  giusto, e solo se può farlo senza bruciare la copertura. Se non trova
  il momento giusto, brucia comunque. Ha già deciso che questa è la sua
  ultima notte come spia.}
\end{secretbox}

\clearpage

% ════════════════════════════════════════════════════════════
% PG 7 — DOTTOR ANSALDO GIUSTINIANI
% @rpg-schema: <https://rpg-schema.org/instances/genova1507/pg-ansaldo-giustiniani>
%   a fitd:Character ;
%   rdfs:label "Dottor Ansaldo Giustiniani" ;
%   fitd:crew <.../crew-fieschi> ;
%   fitd:role "Il Medico di Piazza / L'Intelligence Popolare" ;
%   fitd:action [ fitd:actionName "Study" ; fitd:dots 4 ] ;
%   fitd:action [ fitd:actionName "Attune" ; fitd:dots 3 ] ;
%   fitd:action [ fitd:actionName "Consort" ; fitd:dots 3 ] .
% ════════════════════════════════════════════════════════════

\pgheader[FieschiGreen]{Dottor Ansaldo Giustiniani}{Il Medico di Piazza / L'Intelligence Popolare}{Fieschi}
\pgquote{Curo tutti: nobili, popolani, francesi. Chi paga e chi non può. Questo mi permette di essere ovunque. È un privilegio che non ho mai smesso di sfruttare.}

\section*{Background \& Motivazione}

\begin{rulebox}{}
  \textbf{Background:} Quarantasei anni, medico fisico laureato a Padova, con studio
  in Portoria. Non è un Fieschi — è alleato per vocazione e opportunità. Cura chiunque,
  senza distinzione di casata o nazionalità, il che gli permette di muoversi in ambienti
  che nessun altro PG di questo tavolo può frequentare. Ha curato ufficiali francesi,
  quindi conosce i loro volti, i loro nomi, le loro debolezze fisiche. Ha curato
  popolani dei caruggi, quindi sa quello che Margherita sa — ma con un registro diverso,
  più medico, più freddo.

  \medskip\noindent\textcolor{FieschiGreen}{\textbf{Motivazione:}} L'occupazione
  francese ha interrotto i commerci di farmaci con le spezierie orientali.
  Ha pazienti che muoiono per mancanza di medicine che non riesce a importare.
  Non è una posizione ideologica — è una questione di morti evitabili.
\end{rulebox}

\section*{Attributi \& Azioni}


\begin{attributoblock}[FieschiGreen]{INSIGHT — Mente}
  \azione{Caccia}{Hunt}{1}{Trova persone attraverso i pazienti e i colleghi}
  \azione{Studio}{Study}{4}{Medicina, anatomia, farmacologia, chimica}
  \azione{Osserva}{Survey}{2}{Legge i sintomi delle persone come dei corpi}
  \azione{Ingegno}{Tinker}{3}{Veleni, antidoti, preparati chimici, meccanismi}
\end{attributoblock}


\begin{attributoblock}[FieschiGreen]{PROWESS — Corpo}
  \azione{Finezza}{Finesse}{2}{Mani chirurgiche — precise in ogni contesto}
  \azione{Ombra}{Prowl}{1}{Si muove negli ambienti medici senza farsi notare}
  \azione{Combattimento}{Skirmish}{0}{Non combatte — tratta i feriti di entrambe le parti}
  \azione{Forza bruta}{Wreck}{0}{No}
\end{attributoblock}


\begin{attributoblock}[FieschiGreen]{RESOLVE — Spirito}
      \azione{Percezione}{Attune}{3}{Sente la paura, il dolore, la disperazione come dati clinici}
    \azione{Comando}{Command}{1}{Autorità medica — la gente obbedisce ai medici per paura}
    \azione{Relazioni}{Consort}{3}{Pazienti di ogni ceto, colleghi, speziali, infermieri}
    \azione{Persuasione}{Sway}{2}{Convince con la prognosi — ottimista o catastrofica a seconda}
  \end{attributoblock}

\medskip
\begin{pgclockbox}{Stress \hfill \clock{9}{0}}
  \small\textit{Traumi: $\square$ Freddo \quad $\square$ Duro \quad $\square$ Ossessionato \quad $\square$ Paranoico \quad $\square$ Irrazionale \quad $\square$ Spezzato}
\end{pgclockbox}

\section*{Abilità Speciale}

% @rpg-schema: <.../pg-ansaldo-giustiniani>
%   fitd:specialAbility [ rdfs:label "Il Medico Cura Tutti" ] .
\begin{fieschibox}{Il Medico Cura Tutti}
  Ansaldo ha accesso a qualsiasi ambiente in cui ci siano feriti o malati —
  inclusi presidi francesi, palazzi nobiliari avversari, e zone di conflitto.
  Nessuno ferma un medico durante una battaglia.
  Inoltre, può stabilizzare qualsiasi PG ferito senza tiro, impedendo
  che una Conseguenza Grave diventi permanente.

  \medskip\noindent\textit{Sempre attivo. La stabilizzazione richiede
  che Ansaldo sia fisicamente presente e abbia almeno due minuti di tempo.}
\end{fieschibox}

\section*{Legami}

\begin{table}[h]
\centering\renewcommand{\arraystretch}{1.4}
\begin{tabularx}{\linewidth}{@{}L{2.6cm} X@{}}
  \toprule\rowcolor{FieschiGreen}
  \textcolor{white}{\textbf{PG}} & \textcolor{white}{\textbf{Legame}} \\
  \midrule
  Niccolò & Il conte mi ha chiamato quando è stato ferito in una caduta da cavallo l'anno scorso. Non ha detto a nessuno che era ubriaco. Nemmeno io. \\
  \rowcolor{FieschiGreen!8}
  Margherita & La conosco da vent'anni. Lei porta i malati da me quando non possono pagare. Io li curo. Non è carità — è la clinica più efficiente di Genova. \\
  Frà Tomaso & Il predicatore ha bisogno di qualcuno che gli dica quando si sta ammalando dal troppo parlare. Quello sono io. \\
  \rowcolor{FieschiGreen!8}
  Violante & La spia. Le fornisco una casa sicura quando ne ha bisogno. Non le chiedo perché. \\
  \bottomrule
\end{tabularx}
\end{table}

\section*{Equipaggiamento}
\begin{goldlist}
  \item Borsa medica completa — riconoscibile da qualsiasi soldato di qualsiasi parte
  \item Preparato soporifero in fiala da dieci dosi — «sedativo per procedure dolorose»
  \item Lista dei pazienti attuali nell'amministrazione francese, con diagnosi e orari di visita
  \item Chiave dello studio — e del locale sul retro dove nasconde Violante quando serve
\end{goldlist}


\begin{secretbox}
  \textbf{Segreto:} Il dottor Ansaldo sta curando segretamente il governatore
  Trivulzio per una condizione che, se diventasse pubblica, distruggerebbe
  la sua autorità. La condizione non è mortale ma è umiliante. Trivulzio
  gli ha promesso immunità personale qualunque cosa succeda questa notte —
  ma Ansaldo sa che le promesse di un governatore in fuga valgono poco.
  Ha tenuto un registro clinico completo della condizione di Trivulzio.
  È il documento più pericoloso che esista in questo momento a Genova.

  \medskip\noindent\textcolor{SecretRed}{\textbf{Agenda:}}
  \textit{Garantire che questa notte finisca senza che lui debba usare
  quel registro. Se deve usarlo, lo userà. Ma preferirebbe non farlo —
  non per lealtà a Trivulzio, ma perché usarlo significherebbe smettere
  di essere un medico e diventare qualcos'altro.}
\end{secretbox}

